\documentclass[11pt]{article}

\usepackage{amsmath}
\usepackage{amssymb}
\usepackage{amsthm}

\newtheorem{theorem}{Theorem}
\newtheorem{lemma}[theorem]{Lemma}
\newtheorem{corollary}[theorem]{Corollary}
\newtheorem{proposition}[theorem]{Proposition}
\theoremstyle{definition}
\newtheorem{definition}[theorem]{Definition}
\theoremstyle{remark}
\newtheorem{remark}[theorem]{Remark}

\title{Disproof of conjecture \texttt{00000000159}\\
\large Decompositions of complete graphs of prime order into distinct prime cycles}
\author{earthking11}
\date{12 September 2026}

\begin{document}

\maketitle

\begin{center}
\textbf{Verdict: FALSE}
\end{center}

\section{The conjecture}

\begin{quote}
\textbf{Conjecture.} \emph{For every prime $p \ge 11$, the complete graph
$K_p$ can be decomposed into a union of cycles of pairwise distinct prime
lengths.}
\end{quote}

We disprove the conjecture by exhibiting a single prime $p$ for which it fails,
namely $p = 11$. The obstruction is a plain edge count and requires no
structural graph theory.

\section{The edge-count obstruction}

Throughout, a \emph{cycle} means a simple cycle of a simple graph, in the
standard sense: a closed walk with no repeated vertices other than the starting
and ending vertex. Its length is its number of edges.

\begin{lemma}[Cycle lengths in $K_p$]
\label{lem:length}
Let $C$ be a cycle in $K_p$. Then the length of $C$ equals the number of
distinct vertices it uses, and
\[
3 \le \lvert C \rvert \le p .
\]
\end{lemma}

\begin{proof}
A simple graph has no loops and no parallel edges, so a cycle uses at least
three distinct vertices; hence its length is at least $3$. A cycle in $K_p$
uses pairwise distinct vertices of the $p$-element vertex set, so it uses at
most $p$ vertices. Each edge of the cycle joins two consecutive vertices, and
since the vertices are distinct this is a bijection between the edges of $C$
and its vertices. Therefore the length of $C$ equals its number of distinct
vertices, which lies in $[3, p]$.
\end{proof}

\begin{corollary}[Available lengths for $p = 11$]
\label{cor:lengths}
A cycle of $K_{11}$ whose length is a prime lies in $\{3,5,7,11\}$.
\end{corollary}

\begin{proof}
By Lemma~\ref{lem:length} the length is a prime in $[3,11]$. The primes in that
interval are $3,5,7,11$.
\end{proof}

\begin{lemma}[The available total is at most $26$]
\label{lem:sum}
If $L$ is a set of pairwise distinct primes, each an available cycle length in
$K_{11}$, then
\[
\sum_{q \in L} q \;\le\; 3 + 5 + 7 + 11 \;=\; 26 .
\]
\end{lemma}

\begin{proof}
By Corollary~\ref{cor:lengths} every element of $L$ belongs to
$\{3,5,7,11\}$. Since the elements are pairwise distinct, $L$ is a subset of
this four-element set, so its sum is at most the sum of the whole set, $26$.
\end{proof}

\begin{theorem}
\label{thm:main}
There is no decomposition of $K_{11}$ into cycles of pairwise distinct prime
lengths.
\end{theorem}

\begin{proof}
Let the cycles of a purported decomposition have pairwise distinct prime lengths
$L = \{q_1, \dots, q_r\}$. In either natural reading of ``decomposed into a
union of cycles'':

\begin{itemize}
\item if the decomposition is edge-disjoint, then the cycle lengths sum exactly
      to the number of edges of $K_{11}$, and
\item if the cycles merely cover all edges, then the cycle lengths sum to at
      least the number of edges of $K_{11}$ (the union has at most
      $\sum_i q_i$ edges, counted with multiplicity).
\end{itemize}

Thus in both cases
\[
\sum_{q \in L} q \;\ge\; \lvert E(K_{11}) \rvert \;=\; \binom{11}{2}
\;=\; \frac{11 \cdot 10}{2} \;=\; 55 .
\]
But Lemma~\ref{lem:sum} gives $\sum_{q\in L} q \le 26$, a contradiction.
Therefore no such decomposition exists.
\end{proof}

\begin{corollary}
Conjecture \texttt{00000000159} is false, because it asserts the claim for
\emph{every} prime $p \ge 11$, and it already fails at $p = 11$.
\end{corollary}

\section{The failure is robust and worsens}

The same count obstructs every prime in the conjectured range. Writing
$\operatorname{sop}(p) = \sum_{3 \le q \le p,\ q \text{ prime}} q$ for the total
of all available prime lengths, we need
\[
\operatorname{sop}(p) \;\ge\; \binom{p}{2} \;=\; \frac{p(p-1)}{2}
\]
for a decomposition to be even conceivable. The table below records the first
few primes.

\begin{center}
\begin{tabular}{r|r|r|r}
$p$ & $\binom{p}{2}$ & $\operatorname{sop}(p)$ & verdict \\ \hline
$11$ & $55$  & $26$  & impossible \\
$13$ & $78$  & $39$  & impossible \\
$17$ & $136$ & $56$  & impossible \\
$19$ & $171$ & $75$  & impossible \\
$23$ & $253$ & $98$  & impossible \\
\end{tabular}
\end{center}

The inequality fails for every prime $p$ with $11 \le p \le 2000$; this is
verified exhaustively by \texttt{reproduce.py}. It also fails for all larger
primes, since by the prime number theorem
\[
\operatorname{sop}(p) \sim \frac{p^2}{2 \ln p}
\qquad\text{while}\qquad
\binom{p}{2} = \frac{p^2}{2}\bigl(1 - \tfrac1p\bigr),
\]
so $\operatorname{sop}(p) / \binom{p}{2} \sim 1/\ln p \to 0$. The gap grows
without bound, so the counting obstruction can never be repaired by taking
$p$ larger.

\section{The strongest objection, and its rebuttal}

The only way to escape the obstruction is to change the meaning of ``cycle''.
If ``cycle'' were weakened to mean a \emph{closed trail} --- a closed walk whose
edges are distinct but whose vertices may repeat --- then a closed trail in
$K_p$ could have length much larger than $p$, and the bound
$\lvert C \rvert \le p$ of Lemma~\ref{lem:length} would disappear. With that
weakened reading the edge count alone would no longer obstruct the claim, and
the conjecture could conceivably be true.

This escape is not available. In graph theory a cycle is a simple cycle; a
closed walk with repeated vertices is a closed trail, not a cycle, and the
conjecture says ``cycles''. Under the standard reading Lemma~\ref{lem:length}
applies and the obstruction is decisive. We record the reading we use
explicitly, and we note in passing that the prime $2$ is never an available
length, because every cycle of a simple graph has length at least $3$; this
only makes the available total smaller and the obstruction stronger.

\section{Scope and formalisation}

The LaTeX argument above is complete. The subdirectory \texttt{lean4/}
contains a Lean~4 development, using only \texttt{Std} (no Mathlib), which
formalises precisely the arithmetic core of the obstruction:

\begin{itemize}
\item \texttt{validLengths = [3,5,7,11]} and the fact that every prime in
      $[3,11]$ is a member (a finite check over \texttt{List.range 12});
\item \texttt{validLengths.sum = 26};
\item the central counting lemma: a duplicate-free list all of whose members
      belong to \texttt{validLengths} has sum at most $26$;
\item the corollary that no such list has sum $55$;
\item the final theorem \texttt{conjecture\_00000000159\_false}, which takes
      such a list together with the hypothesis that its sum equals the edge
      count $11 \cdot 10 / 2$ and derives \texttt{False}.
\end{itemize}

The graph-theoretic step --- that the cycle lengths of a cover of $K_{11}$ form
such a list and must total at least $\binom{11}{2}$ --- is not formalised; see
\texttt{lean4/README.md} for the exact boundary of the formalisation. The
formalised statement is not vacuous: it genuinely asserts that no list of
pairwise distinct available prime lengths can reach the required total.

\section{Files and reproduction}

\begin{center}
\begin{tabular}{l|l}
file & purpose \\ \hline
\texttt{main.tex} & this disproof \\
\texttt{build/main.pdf} & compiled PDF (produced by \texttt{tectonic main.tex}) \\
\texttt{reproduce.py} & verifies the count for all primes up to $2000$ \\
\texttt{lean4/} & Lean~4 formalisation, core library only, no \texttt{sorry} \\
\end{tabular}
\end{center}

Reproduce the numerical check with pure Python~3 (standard library only):

\begin{verbatim}
python3 reproduce.py
\end{verbatim}

It prints the table for $p = 11, 13, 17, 19, 23$, lists the four available
lengths at $p = 11$ with $3+5+7+11 = 26 < 55$, asserts that the bound fails for
every prime in $[11,2000]$, and ends with a \texttt{PASS} summary.

\section*{Status against the submission rules}

Rule 3 requires a LaTeX source, a compiled PDF, and a Lean~4 project. All three
are present.

\begin{itemize}
\item \textbf{LaTeX source:} \texttt{main.tex}.
\item \textbf{Compiled PDF:} \texttt{build/main.pdf}, produced by
      \texttt{tectonic main.tex}. The source needs only the standard
      \texttt{article}/\texttt{amsmath}/\texttt{amssymb}/\texttt{amsthm}
      packages. (It was not regenerated in the offline environment in which
      this submission was assembled, since the TeX resource bundle could not
      be fetched there.)
\item \textbf{Lean~4 project:} \texttt{lean4/}, built with \texttt{lake build},
      using core Lean plus \texttt{Std} only (no Mathlib), with no
      \texttt{sorry}. The axiom footprint is machine-checkable via
      \texttt{lake env lean Check.lean}.
\end{itemize}

\end{document}
